\documentclass[11pt]{article}

% ---------------------------------------------------------------------------
%  L8 — Below the onset: the diagonal law's error term.
%  Category L (entirely machine-written); see paper/README.md and
%  docs/publication-split.md.
%
%  NOVELTY: pass run 2026-08-18, no collision found -- but a weak negative, see
%  the Novelty section.  docs/priority-passes-2026-08-18.md.  Claims none.
%
%  Source material, all in this repository:
%    docs/proofs/diagonal-law.md              the law, its onset, its proof
%    results/diagonal-law-below-onset.md      how the law fails below onset
%    results/onset-defect-law.md              rate 9, exponent j-3/2, amplitude
%    results/onset-defect-depth1-closed.md    the quartic, and every j=1 constant
%    results/onset-defect-depths234.md        depths 2,3,4 closed the same way
%    results/ridgeline-depth-amplitudes.md    alpha = 50/81, and the limits ledger
%    results/onset-defect-crossover.md        the near-onset layer, downgraded
%    results/allpairs-kernel.md               the gap walk
% ---------------------------------------------------------------------------

% ---------------------------------------------------------------------------
%  paper/shared/preamble.tex — packages and macros common to every manuscript
%  in this directory.  Factored out of paper/polyplets-report.tex and
%  paper/technical-report.tex, which between them had two copies of the same
%  twenty lines.
%
%  technical-report.tex does NOT \input this file and is not to be edited to do
%  so: it is jasonp's prose and is read-only to the machine (paper/README.md,
%  "Who may edit what").  The duplication there is deliberate.
%
%  Assumes \documentclass[11pt]{article} has already been issued.
% ---------------------------------------------------------------------------
\usepackage[margin=1in]{geometry}
\usepackage{amsmath,amssymb}
\usepackage{amsthm}
\usepackage{booktabs}
\usepackage{graphicx}
\usepackage{numprint}
\usepackage{microtype}
\usepackage[table]{xcolor}
\usepackage[hidelinks]{hyperref}
\usepackage{flafter}   % a float never appears before the text that introduces it
\usepackage{float}     % [H] pins tables/figures exactly where written
\usepackage[font=small,labelfont=bf,labelsep=period]{caption}
\setlength{\intextsep}{16pt plus 3pt minus 3pt}

\npthousandsep{,}
\npfourdigitsep
\newcommand{\num}[1]{\numprint{#1}}
\newcommand{\g}{\cellcolor{gray!15}}

% Theorem environments.  One counter shared across kinds, so that a reference
% like "Theorem 4" is unambiguous in a paper that also has lemmas and
% propositions.
\theoremstyle{plain}
\newtheorem{theorem}{Theorem}
\newtheorem{proposition}[theorem]{Proposition}
\newtheorem{lemma}[theorem]{Lemma}
\newtheorem{corollary}[theorem]{Corollary}
\theoremstyle{definition}
\newtheorem{definition}[theorem]{Definition}
\newtheorem{example}[theorem]{Example}
\theoremstyle{remark}
\newtheorem{remark}[theorem]{Remark}
\newtheorem{conjecture}[theorem]{Conjecture}
\newtheorem{openproblem}[theorem]{Open Problem}

% Repo-path citations.  Every one of these must resolve in the tree that ships
% with the paper; `make gate-citations` checks the markdown, and
% paper/verify_claims.py checks the manuscripts.
\newcommand{\repo}[1]{\texttt{#1}}

% OEIS links.
\newcommand{\oeis}[1]{\href{https://oeis.org/#1}{#1}}

% Growth constants, so a rename is one edit.
\newcommand{\lam}{\lambda}
\newcommand{\muH}{\mu_H}

% ---------------------------------------------------------------------------
%  paper/shared/disclosure.tex — the two authorship disclosure blocks.
%
%  docs/publication-split.md §1 fixes the rule these implement:
%
%    "Under no circumstances will there be any ambiguity about whether a paper
%     is wholly my work (none), merely my prose (some) or entirely LLM-authored
%     (some)."   — jasonp, 2026-08-07
%
%  There are exactly two blocks and they live in exactly one file, so that a
%  paper cannot quietly soften its own disclosure.  The fixed sentences take no
%  arguments.  The ONE thing a paper supplies is the per-result verification
%  ledger, because that is the part that differs honestly between papers — and
%  §1 requires it to be per result, not in aggregate.
%
%  Do not add a \Pdisclosure or \Ldisclosure variant with an optional softening
%  argument.  If a paper does not fit one of these two blocks, the paper is
%  wrong, not the block.
% ---------------------------------------------------------------------------

\newsavebox{\disclosurebox}

\newenvironment{disclosureframe}
  {\begin{center}\begin{lrbox}{\disclosurebox}\begin{minipage}{0.93\textwidth}\small}
  {\end{minipage}\end{lrbox}\fbox{\usebox{\disclosurebox}}\end{center}}

% --- P: jasonp's prose -------------------------------------------------------
% Byline: Jason Parker, alone.  Argument: the per-result ledger of what the
% machine did.
\newcommand{\Pdisclosure}[1]{%
\begin{disclosureframe}
\textbf{Authorship.}\enspace Every sentence of this paper was written by its
named author. He has read every proof it states and believes each one, and he
is responsible for the correctness of what it claims.

\smallskip
\textbf{Machine contribution.}\enspace #1
\end{disclosureframe}}

% --- L: entirely machine-written --------------------------------------------
% Argument: the per-result verification ledger — for each result, what warrants
% it (Lean kernel, exact certificate, or labelled measurement) and how much of
% it a human has checked.
\newcommand{\Ldisclosure}[1]{%
\begin{disclosureframe}
\textbf{Authorship.}\enspace This document was written end to end by a large
language model. That includes its mathematics, its prose, its tables and its
\LaTeX{}. No sentence of it was written by a human. The mathematics it reports
was likewise derived by machine.

\smallskip
\textbf{Human role.}\enspace The person who directed this work chose what to
compute, chose what to publish, and is responsible for the decision to release
this document. Except where the ledger below says otherwise, that person has
not personally checked the arguments here, and this disclosure is not to be
read as an endorsement of them.

\smallskip
\textbf{What warrants each result.}\enspace #1
\end{disclosureframe}}

% --- Draft banner ------------------------------------------------------------
% For an L-paper whose verification ledger is not yet settled.  Loud on purpose:
% an L-paper without a truthful ledger must not be mistaken for a finished one.
%
% \firstdraftbanner is the standard wording, which five papers previously
% carried character-for-character in their own preambles.  It lives here for the
% same reason the disclosure blocks do: identical text in seven places drifts,
% and a paper must not be able to soften it by editing its own copy.  The
% optional argument is for a gate specific to one paper, appended verbatim; a
% paper needing different wording (L6, compute-gated) calls \draftbanner direct.
\newcommand{\firstdraftbanner}[1][]{%
\draftbanner{This is a first draft. The verification ledger below records what
warrants each result \emph{mechanically}; the column that records what a human
has read is empty, deliberately, because nobody has read it yet. Do not
circulate until that column says something true.#1}}

\newcommand{\draftbanner}[1]{%
\begin{center}
\fcolorbox{red!60!black}{yellow!15}{%
  \begin{minipage}{0.93\textwidth}\small
  \textbf{DRAFT — NOT FOR CIRCULATION.}\enspace #1
  \end{minipage}}
\end{center}}


\newcommand{\Z}{\mathbb{Z}}
\newcommand{\Q}{\mathbb{Q}}
\newcommand{\R}{\mathbb{R}}

\title{\textbf{Below the onset}\\[4pt]
\large The error term of a diagonal law, depth by depth}
\author{[byline: jasonp's call --- \repo{docs/publication-split.md} \S1]}
\date{August 2026}

\begin{document}
\maketitle

\firstdraftbanner[ One gate is specific to this paper: its priority pass
(\S\ref{sec:novelty}) is a weak negative, because the object it studies is
defined relative to a law of our own. Nothing here is claimed to be new.]

\Ldisclosure{%
\emph{The diagonal law itself, its degree, its onset and its sharpness} ---
proved elsewhere (companion paper L1), used here as stated.
\emph{The exact frame of \S\ref{sec:frame}} --- proved: the entire below-onset
content of column $k$ is the $k+1$ coefficients of one correction polynomial,
which is a rearrangement of the chain identity and nothing more.
\emph{The depth-$1$ closure of \S\ref{sec:depth1}} --- the defect is computed
exactly to $k = 200$ from independently enumerated cluster families, and the
quartic $\Phi$ was fitted on $57$ series orders and holds on $144$ orders of
pure holdout, with a mod-$p$ scan showing nothing smaller works. The constants
$9$, $-1/2$, $\sqrt6/(27\sqrt\pi)$ and $a \in \Q(\sqrt3)$ are then
\emph{derived} from $\Phi$ by singularity analysis, not fitted.
\emph{Depths $2,3,4$ (\S\ref{sec:depths})} --- exact identities over
independently enumerated families, agreeing with the banked triangle at
$19+18+17+16$ cells, with a RED control that fires.
\emph{$\theta_j = j - 3/2$ and rate $9$ at $j \le 7$ (\S\ref{sec:allj})} ---
labelled measurement, bias-calibrated against controls whose correction series
do \emph{not} terminate, and the calibration is self-consistent rather than
independent: it removes an artefact, it does not prove the law.
\emph{$\alpha = 50/81$ (\S\ref{sec:ridge})} --- derived, by a finite local
enumeration ($\Sigma = 450$ out of $450$ excursions), \textbf{conditional on
four named analytic assumptions} which are listed in full in
\S\ref{sec:assumptions} and are not proved anywhere.
\emph{The amplitude family at $j \ge 5$} --- a \textbf{prediction}. No exact
$D_5$ beyond $k \le 19$ exists.
\emph{The near-onset scaling of \S\ref{sec:crossover}} --- measurement,
downgraded on adversarial review and reported at the strength that survived it.
\emph{The square-lattice check of \S\ref{sec:square}} --- \textbf{external}:
the counts are other people's. At $k \le 6$ the below-onset fit equals the
classical one and reproduces the cells it was denied, ten of them at depth $2$,
with a RED control at each depth that fires. Its limit is stated in place ---
the square defects are measured against a classically pinned $P_k$, so this
validates the identity and is not an independent route to $D_j$ off the king
lattice.
\emph{Novelty} --- pass run 2026-08-18, no collision found, and
\S\ref{sec:novelty} says why that is weaker evidence than it looks. The
square-lattice material of \S\ref{sec:square} postdates that pass and has not
been swept.
\emph{Human verification:} none, as of this draft.}

\begin{abstract}
\noindent
On the king lattice, the number $T(n,H)$ of fixed animals of $n$ cells with
bounding-box height exactly $H$ satisfies a diagonal law: for $k$ fixed and
$n \ge 2k+1$,
\[
  T(n, n-k) \;=\; P_k(n)\cdot 3^{\,n-1-3k},
\]
with $P_k$ a polynomial of degree $k$. The onset $n \ge 2k+1$ is sharp. This
paper is about what happens below it, where the law is false.

It is false by a vanishing relative amount, and the error has structure that no
algebraic search finds because the structure is analytic. Writing the defect at
depth $j = 2k+1-n$ as $D_j(k)$, we report: an exact frame in which the entire
below-onset content of a column is $k+1$ integers; a closed form at depth $1$,
where the generating function is algebraic, an irreducible quartic, from
which the rate $9$, the exponent $-1/2$, the amplitude $\sqrt6/(27\sqrt\pi)$ and
the next coefficient $a = 3293/92928 - 3251\sqrt3/185856 \in \Q(\sqrt3)$ all
follow by singularity analysis; the same closure at depths $2$, $3$ and $4$; and
the measurement $\theta_j = j - 3/2$ with rate $9$ at every depth $j \le 7$.

The depth direction then collapses. The seven measured amplitude ratios are not
seven numbers but one: the velocity $\alpha = 50/81$ at which the defect's
branch point moves when a cluster is allowed one unit of excess, computed exactly
by a finite local count. That resolves an ambiguity the measurements alone could
not --- at $j = 5$ the value is $35/8$, and the competing $118/27$, which fits
the numbers better, is an artefact of the normalisation in which they were read.

The derivation of $\alpha$ rests on four analytic assumptions which we state
rather than prove.

Finally, the construction is checked outside this project. On the square
lattice, where the counts are published by other people, a below-onset cell
corrected by its defect gives the same polynomial as the classical fit and
reproduces the tall cells it was denied, at both depths tested.
\end{abstract}

\tableofcontents

\section{Introduction}
\label{sec:intro}

A diagonal law is an exact statement with an onset. Companion paper L1 proves
that for a row-local lattice with drift $b$, the diagonal $T(H+k, H)$ equals
$q_k(H)b^H$ with $\deg q_k \le k$, for $H$ past an explicit threshold, and that
the threshold is sharp: at one row below it the formula is wrong.

Sharpness is usually where a result stops. This paper starts there.
The king-lattice instance, in the indexing used throughout,
\begin{equation}
  T(n, n-k) \;=\; P_k(n)\cdot 3^{\,n-1-3k}, \qquad n \ge 2k+1,
  \label{eq:law}
\end{equation}
is exact above the onset line and false below it, and the region below the onset
line is not a boundary case --- it is the middle of the triangle, $H \le n/2$.
Define the \emph{depth}
\[
  j \;=\; 2k+1-n \;\ge\; 1
\]
and the \emph{defect}
\[
  D_j(k) \;=\; T(n,n-k) - P_k(n)\cdot 3^{\,n-1-3k}, \qquad n = 2k+1-j .
\]
Depth $1$ is the first row the law misses, depth $2$ the next, and so on.

Two facts frame everything below. First, the law is asymptotically exact where
it is false: the number of correct leading digits of $T$ grows with $k$ at every
fixed depth. Second, the defect has no algebraic structure that a search finds
directly --- C-finite to order $8$ and P-finite to $(r,d) = (4,4)$ both come back
null. The reason, established here, is that those searches were looking in the
wrong category and at the wrong scale: $D_1$ \emph{is} P-finite, at
$(r,d) = (35,4)$, and the structure that explains the whole family is a moving
branch point rather than a recurrence.

\paragraph{What this paper is not.} It is not a proof of the defect law. Rate
$9$ and exponent $j - 3/2$ are proved at $j = 1$ and measured at $j \le 7$;
the amplitude family is derived from an identity plus four analytic assumptions
that are stated in \S\ref{sec:assumptions} and proved nowhere. The ledger on
page 1 says which is which, claim by claim, and \S\ref{sec:novelty} says why the
literature search that was run is weak evidence either way.

\section{The exact frame: a column is \texorpdfstring{$k+1$}{k+1} integers}
\label{sec:frame}

The chain identity behind \eqref{eq:law} gives, per column $k = n - H$,
\begin{equation}
  [y^k]F \;=\; \frac{R_k(z)}{(1-3z)^{\,k+1}}, \qquad R_k \in \Z[z], \quad
  \deg R_k \le 2k+1,
  \label{eq:column}
\end{equation}
and the step that produces the law splits off a correction polynomial $D(z)$ of
degree at most $k$. Everything below the onset is that polynomial:
\begin{equation}
  D_j(k) \;=\; [z^{\,k+1-j}]\,D(z).
  \label{eq:defectcoeff}
\end{equation}
So the below-onset part of a column is $k+1$ integers rather than an infinite
tail to be estimated, and the depth-$j$ defect is one of them. In
particular the top one is a single coefficient of $R_k$:
\begin{equation}
  D_1(k) \;=\; \frac{\operatorname{lead}(R_k)}{(-3)^{\,k+1}}.
  \label{eq:d1lead}
\end{equation}

The rest is computable because of this. A defect defined as a difference of two
large numbers would inherit their size; \eqref{eq:defectcoeff} makes it an object
in its own right, with its own generating function.

\begin{remark}[what the $k+1$ integers do not buy]
Two readings of \eqref{eq:defectcoeff} are natural and both are wrong, so they
are recorded here rather than left for a reader to try.

\emph{They are not surplus equations.} A column below the onset holds $k+1$
integers, which looks like $k+1$ constraints on a level that needs two. It is
not: each below-onset cell brings one equation \emph{and} one unknown $D_j(k)$,
so solving a whole column is a reformulation rather than extra redundancy. The
honest accounting is
\[
  \text{surplus} \;=\; \bigl(\text{cells whose depth has a known } D_j\bigr) - 2,
\]
and a cell at a depth whose defect is unknown contributes nothing.

\emph{And they are not enough to fit a recurrence.} Nothing above rules out
$D_j$ being P-finite in $k$, and an empirical search is the obvious thing to
try. It is priced out rather than untried: fitting a recurrence of the orders in
question needs on the order of $180$ values of $k$, and about $15$ exist.
\end{remark}

\section{Depth one closes}
\label{sec:depth1}

\subsection{The defect without the triangle}

The obstacle to computing $D_1(k)$ far was that $P_k$ was known only to
$k \le 19$. Equation~\eqref{eq:d1lead} removes the obstacle: the leading
coefficient of $R_k$ assembles from the all-pairs cluster families alone, the
gap walk, with no triangle data and no wired $P_k$ anywhere. Depth one is
therefore exact to $k = 200$ in about $40$ seconds, against a previous ceiling
of $19$.

\subsection{The generating function is algebraic}

With $200$ exact terms in hand, the generating function
$W(x) = \sum_k D_1(k)x^k$ satisfies an irreducible quartic $\Phi(x, W) = 0$.
The fit used $57$ series orders; the quartic then holds on $144$ further orders
that the fit never saw, and a modular scan shows no smaller box works --- no
cubic, no quadratic, nothing rational, within the degrees the series supports.

The earlier null results missed that the object is algebraic.
$D_1$ is also P-finite, at $(r,d) = (35,4)$; the verdicts of ``not P-finite'' at
envelopes $(4,4)$ and $(6,8)$ were correct as scoped and are explained by the
size of the true one.

\subsection{Every depth-one constant, derived}

From $\Phi$, singularity analysis~\cite[Ch.~VI--VII]{flajoletSedgewick2009}
gives the constants that the campaign had been measuring; $\Phi$ is algebraic
with a square-root branch point, which is the case
\cite[Thm.~VII.8]{flajoletSedgewick2009} covers, and the $k^{-1/2}$ and the
$1/k$ correction are read off the Puiseux expansion there:
\[
  D_1(k) \;=\; C_1\,9^k\,k^{-1/2}\Bigl(1 + \frac{a}{k} + O(k^{-2})\Bigr),
  \qquad
  C_1 = \frac{\sqrt6}{27\sqrt\pi},
\]
\[
  a \;=\; \frac{3293}{92928} - \frac{3251\sqrt3}{185856}
    \;=\; 0.005138939956706\dots
\]
The rate $9$ and the exponent $-1/2$ are branch data of $\Phi$. The amplitude
matches the measured $0.05118432$ within a control-backed bar of $3\times10^{-7}$
relative, and the recognition is unique: over all $\sqrt m/(n\sqrt\pi)$ with
$m \le 200$, $n \le 400$, exactly one value lies within $10^{-6}$ of the
measurement.

The coefficient $a$ is the useful one, because it explains a failure. It lies in
$\Q(\sqrt3)$, which is why every attempt to recognise it at $k \le 19$ over the
rationals, or over products of $\sqrt6$ and $\pi$, had to fail. Nothing was wrong with
those searches except the field they searched in.

\section{Depths two, three and four}
\label{sec:depths}

The same closure runs at $j = 2, 3, 4$: an exact identity over independently
enumerated weight families, with no triangle data and no wired $P_k$ in the
derivation. Grade every cluster family by its excess
$e = \sum_i (s_i - 2)$. Then the coefficient $[z^{k+1-j}]$ of the correction
polynomial draws only on families of total excess $\le j-1$, assembled through
the chain identity with binomial $(1-3z)$ corrections.

The grading by excess is the substantive step. It is not ``up to $j-1$ rows of
size $3$'' --- a first reading of the depth-$1$ case suggests that, and it is
wrong:
depth $3$ needs the one-$4$-row family, and depth $4$ needs one-$5$-row, $4+3$,
and $3+3+3$.

\begin{table}[H]
\centering\small
\begin{tabular}{l l l}
\toprule
depth & banked cells matched exactly & range \\
\midrule
$j = 1$ & $19$ & $k = 1..19$ \\
$j = 2$ & $18$ & $k = 2..19$ \\
$j = 3$ & $17$ & $k = 3..19$ \\
$j = 4$ & $16$ & $k = 4..19$ \\
\bottomrule
\end{tabular}
\caption{The gate assembles its banked side from the triangle and the wired
$P_k$, which the candidate identities never touch, so the two sides share no
machinery. A RED self-test fires.}
\label{tab:depths}
\end{table}

\section{All depths, measured}
\label{sec:allj}

Beyond $j = 4$ there is no closed form and there is measurement. Write
$D_j(k) \sim C_j\,9^k\,k^{\theta_j}$.

\begin{table}[H]
\centering\small
\begin{tabular}{r l l l}
\toprule
$j$ & $\theta_j$ & $j - 3/2$ & rate \\
\midrule
1 & $-0.5000 \pm 0.0010$ & $-0.5$ & $8.999998 \pm 2.2\times10^{-5}$ \\
2 & $+0.5000 \pm 0.0025$ & $+0.5$ & $9.000000 \pm 2.1\times10^{-3}$ \\
3 & $+1.5002 \pm 0.0065$ & $+1.5$ & $8.9988 \pm 5.6\times10^{-3}$ \\
4 & $+2.5006 \pm 0.0110$ & $+2.5$ & $8.9945 \pm 8.9\times10^{-3}$ \\
5 & $+3.4974 \pm 0.0130$ & $+3.5$ & $9.0395 \pm 5.1\times10^{-2}$ \\
6 & $+4.4999 \pm 0.0090$ & $+4.5$ & $9.0025 \pm 7.2\times10^{-2}$ \\
7 & $+5.4983 \pm 0.0385$ & $+5.5$ & --- \\
\bottomrule
\end{tabular}
\caption{$\theta_j = j - 3/2$ at every depth measured, and rate $9$ inside the
bar at every depth. Bias-calibrated; see below.}
\label{tab:thetas}
\end{table}

\paragraph{The estimator is biased and the bias is the story.} The naive
estimator $L_k = k(D_k/D_{k-1}/9 - 1) \to \theta$ carries a systematic of up to
$0.02$ in $\theta$ and $0.03$ in the rate when run against a control whose
correction series does not terminate. An earlier reading of this data reported
$\theta$ drifting low with $j$ and the rate falling away from $9$: $8.98$,
$8.91$, $8.77$ at $j = 3,4,5$. Both drifts are reproduced exactly by a control
whose $\theta$ and rate are right by construction. They were the estimator.

The correction is not free of assumption, and the ledger says so: the matched
control is built assuming $\theta_j = j - 3/2$, so Table~\ref{tab:thetas} is a
self-consistency test. It removes ``$\theta$ drifts low'' as evidence against the
law; it is not independent evidence for it.

\section{The depth direction is one constant}
\label{sec:ridge}

\subsection{Depth is a coupling, not a new object}

Under $Y = yz$ and $t = 1/z$, the proved chain identity collapses to
\begin{equation}
  G(Y,t) \;=\; \widehat P(Y,t) + \frac{\widehat B(Y,t)^2}{\,t - 3 - \widehat S(Y,t)\,},
  \qquad
  D_j(k) = [Y^k t^{\,j-1}]\,G ,
  \label{eq:master}
\end{equation}
valid for $j - 1 \le k$, where $\widehat S, \widehat B, \widehat P$ are the
excess-graded aggregates of the cluster families and $t$ marks excess. The
grading is legitimate because a family of excess $e$ carries $k \ge e+1$, so
$t^e$ always arrives with at least $Y^{e+1}$.

Equation~\eqref{eq:master} reproduces the depth-$1$ through depth-$4$ series at
all $74$ exact cells with $k \le 19$, and perturbing one excess-$1$ weight by
$+1$ breaks it.

The identity's use is not brevity. It identifies the depth variable as a marker
for cluster excess, so depth is a perturbation of the row transfer of the gap
walk --- a coupling constant rather than a new object to be summed.

\subsection{A moving branch point}

In the normalisation $D_j(k) \sim C_j 9^k k^{j-3/2}$ the tower has one scaling
variable $\tau = t/(1-9Y)$, and the family of amplitude ratios becomes
\[
  \frac{C_{M+1}}{C_1} \;=\; \frac{\alpha^M}{M!},
\]
a pure exponential rather than a family of binomials. An exponential enhancement
per unit chain length is exactly a square-root branch point whose \emph{location}
moves linearly with the excess marker,
\[
  1 - 9Y_c(t) \;=\; \alpha t + O(t^2),
\]
with the exponent pinned at $-1/2$. So $\alpha$ is a local quantity: the
first-order shift of the zero-momentum eigenvalue of the gap walk's row transfer
under one $3$-cell row. That is a finite count, and it is
\[
  \alpha \;=\; \frac{\Sigma}{729} \;=\; \frac{450}{729} \;=\; \frac{50}{81}.
\]

\subsection{What the measurements could not decide}

In the earlier normalisation the amplitude ratios read
$1, 1, 3/2, 5/2, 4.370843, 7.866341, 14.441622$ against the candidate
$\binom{2j-2}{j-1}/2^{j-1} = 1, 1, 3/2, 5/2, 35/8, 63/8, 231/16$. At $j = 3$ and
$j = 4$ the candidate is the only rational with denominator $\le 32$ inside the
measurement bar. At $j = 5$ it is not: eight simple rationals fit, and one of
them, $118/27 = 4.370370$, is \emph{closer} to the measurement than
$35/8 = 4.375$.

The derivation settles it. With $C_{M+1}/C_1 = \alpha^M/M!$ exactly and
$\alpha = 50/81$, the value at $j = 5$ is $35/8$. The measurement returns
$118/27$ when it is read in a normalisation that hides the exponential, so that
number is an artefact of the reading rather than a competing constant.

This is the paper's one instance of a derivation overturning a better numerical
fit. The direction of the evidence is the part to keep in view: the $j = 5$
number alone favours the wrong answer, and only the mechanism behind $\alpha$
selects the right one.

\subsection{The four assumptions}
\label{sec:assumptions}

The derivation of $\alpha$ is conditional. The conditions, stated as they stand
and not softened:

\begin{enumerate}
\item \emph{The exponent stays exactly $-1/2$ for small $t > 0$.} The dispersion
  has a nondegenerate stationary point at $u = 1$, with $(\log K)''(1) = 4/3$
  symbolically, and analytic perturbations of a nondegenerate minimum stay
  nondegenerate. What is missing is control of the step from the dispersion to
  the assembled $G$.
\item \emph{$Y_c(t)$ is analytic at $t = 0$.} Each order in $t$ adds finitely
  many excess families, so the vertex expansion is order-by-order finite.
  Convergence is not proved.
\item \emph{The cancellation that keeps $1/9$ dominant persists at every order in
  $t$.} Rate $9$ is measured at every depth $j \le 6$ and derived at $j = 1$
  only.
\item \emph{The perturbative bookkeeping behind $\alpha = \Sigma/729$} --- a flat
  zero-momentum projection inside the delocalised sector. This was the weakest
  link and is now much less weak: the excursion lands in the well-behaved sector
  $450$ times out of $450$ and in the badly-behaved one never, so the vertex acts
  where the transfer is a translation-invariant walk of mass $9$ and the flat
  mode is the critical mode. What remains unproved is the passage from
  ``the zero-momentum mass shifts by $Y^2t\,\Sigma/9$'' to ``the branch point of
  $G$ moves by $\alpha t$ with its exponent intact'', which is assumption 1 in
  other clothes.
\end{enumerate}

Two further items are verified rather than derived. $\alpha_M$ is constant over
$M = 1,2,3$ within calibrated bias --- sharp at $M=1$ ($1.1\times10^{-9}$),
useful at $M=2$, and at $M=3$ constraining almost nothing on its own. And the
law at $j = 5,6,7$ is a prediction: no exact $D_5$ beyond $k \le 19$ exists here
or anywhere.

\section{The near-onset layer}
\label{sec:crossover}

Close to the onset line a resummation of the family,
\[
  D_j(k) \;\approx\; \frac{\sqrt6}{27}\cdot\frac{\binom{2N}{N}}{4^N}
  \binom{N}{k}9^k\Bigl(\frac{50}{81}\Bigr)^{j-1}, \qquad N = k+j-1,
\]
is accurate and degrades smoothly with depth, which implies a scaling variable.
Fitting the residual surface over $168$ banked below-onset cells finds a boundary
layer of width $j \sim k^{0.4}$.

That exponent, and the crossover curve read off it, were downgraded on
adversarial review and are reported here as measurement with the weaker status
they came back with; the family exclusions from the same analysis stand. The
resummation is a near-onset statement and not a uniform description of the
below-onset region.

\section{An external check: the same construction on the square lattice}
\label{sec:square}

Everything above is internal to one lattice and, in the end, to one project.
Every second source in this repository re-counts the same objects with different
code, so a shared \emph{conceptual} error would survive all of them. The square
lattice is where that stops being true, because its counts are published by
other people.

There the drift set has $b = |D| = 1$, so the $3^{\,n-1-3k}$ factor is $1$ and
the law is a plain polynomial $T_{\mathrm{sq}}(n, n-k) = P_k(n)$ of degree $k$
for $n \ge 2k+1$. The depth-$1$ defect, measured ab initio on the banked square
bounding-box triangle, is
\[
  D_1(k) \;=\; (-1)^{\,k+1},
  \qquad k = 1,\dots,6 .
\]

\paragraph{Depth one.} Fit each $P_k$ twice: classically, from the $k+1$
in-onset cells at $n = 2k+1, \dots, 3k+1$; and by dropping the \emph{tallest} of
those and using the depth-$1$ cell at $n = 2k$ in its place. At every level
$k = 1,\dots,6$ the two fits are equal as polynomials, and the second reproduces
the tall cell it was denied. A RED control confirms that the wrong defect sign
breaks the fit.

\paragraph{Depth two.} Dropping the two tallest in-onset cells and using
$n = 2k-1$ and $n = 2k$ instead gives exactly $k+1$ points for a degree-$k$
polynomial, with the two tallest cells never seen. At $k = 2,\dots,6$ the fit
again agrees with the classical one and predicts both withheld cells exactly:
ten held-out cells, ten exact hits, in exact rational arithmetic. A $D_2$
corrupted by $1$ breaks the fit at every $k$ --- without that second control the
whole table would be consistent with $D_2$ being irrelevant.

\paragraph{And $D_2$ has no law in five levels.} The measured values are
\[
  D_2(k) \;=\; -4,\; 8,\; -3,\; 10,\; -1
  \qquad (k = 2,\dots,6).
\]
The sign alternates as $(-1)^{k+1}$, like $D_1$. The magnitudes $4, 8, 3, 10, 1$
are not monotone, do not separate by parity, and support no fit worth writing
down on three and two points. They are recorded here as five measurements and
not as a law.

\paragraph{The limit, which is the important part.} These square defects are
measured \emph{against} a $P_k$ that was pinned classically. That is exactly
what is needed to validate the identity --- a cell below the onset, corrected,
is a valid equation for the same polynomial --- but it is not an independent
route to $D_j$ on the square lattice, and so it does not by itself make a square
$P_k$ cheaper. On the king lattice the corresponding defects come from the
family enumeration of \S\ref{sec:depth1} and \S\ref{sec:depths}, which never
sees $P_k$. Making that enumeration lattice-parametric is a build and not a
substitution: its row transfer has king adjacency baked in, and the
lattice-parametric machinery this project has elsewhere covers the
\emph{above}-onset assembly only.

\section{Novelty}
\label{sec:novelty}

A pass ran on 2026-08-18 and found nothing close
(\repo{docs/priority-passes-2026-08-18.md}): searches over corrections and error
terms to fixed-height and bounding-box animal formulas, asymptotics of defects
below a formula's threshold, and algebraic generating functions for the leading
coefficients of such families.

That negative is \emph{weak}, for a stateable reason. The object here is
defined relative to a diagonal law of our own, so a collision would have to be
with work that had the law first, which companion paper L1's searches cover, and
they came back clean. A pass over an object nobody else has defined
tests very little. Nothing here is called new.

\section{Open problems}

\begin{openproblem}[prove the defect law]
Rate $9$ and exponent $j - 3/2$ are derived at $j = 1$ from the quartic and
measured at $j \le 7$. A proof at general $j$ needs the four assumptions of
\S\ref{sec:assumptions}, of which item 4 is self-contained: redo the first-order
perturbation with the correct left and right critical vectors of the transfer and
confirm that $\Sigma/729$ falls out.
\end{openproblem}

\begin{openproblem}[make the sharpness argument unconditional and effective]
The amplitude $C_1 = \sqrt6/(27\sqrt\pi)$ is positive, so $D_1(k) \ne 0$ for all
sufficiently large $k$ --- and $D_1(k) \ne 0$ is the non-cancellation
statement that makes the diagonal law's onset sharp, which companion paper L1
records as open. Two gaps stand between that remark and a theorem: the quartic
$\Phi$ is fitted and held on $144$ holdout orders rather than proved, and
``sufficiently large'' is not effective. Closing either is worth more than
another depth.
\end{openproblem}

\begin{openproblem}[an exact $D_5$]
Every statement about $j \ge 5$ in this paper is a prediction, because the exact
data stops at $k \le 19$ there. Extending the excess-graded family enumeration by
one unit of excess would test the $35/8$ against something other than its own
derivation.
\end{openproblem}

\begin{openproblem}[the algebraic tower]
$D_1$ is algebraic. Is $D_2$? The master identity gives each depth as a
coefficient of one bivariate object, so the natural question is whether $G$ is
algebraic in $Y$ at every fixed order in $t$, which the depth-$1$ quartic
establishes only at $t^0$.
\end{openproblem}

\section{Reproduction}
\label{sec:repro}

\begin{table}[H]
\centering\small
\begin{tabular}{@{}l R{0.56\textwidth}@{}}
\toprule
claim & check \\
\midrule
the law, its onset and sharpness & companion paper L1; \repo{docs/proofs/diagonal-law.md} \\
the law is asymptotically exact below onset & \repo{experiments/law\_below\_onset.py} \\
the exact frame, and $D_j$ as a coefficient & \repo{experiments/defect\_below\_onset.py} \\
depth 1 to $k = 200$ & \repo{experiments/depth1\_gap\_walk.py} \\
the quartic $\Phi$ and its holdouts & \repo{experiments/depth1\_minpoly.py} \\
the derived constants $C_1$, $a$ & \repo{experiments/depth1\_asymptotics.py} \\
$D_1$ P-finite at $(35,4)$ & \repo{experiments/depth1\_recurrence.py} \\
depths 2--4, and the gate & \repo{experiments/severance\_w3\_depths.py}, \repo{experiments/severance\_w3\_gate.py} \\
$\theta_j$, the rate, and the controls & \repo{experiments/defect\_nine\_exponent.py}, \repo{experiments/defect\_controls.py} \\
the amplitude and its family & \repo{experiments/defect\_amplitude.py}, \repo{experiments/defect\_amplitude\_family.py} \\
the master identity and $\alpha$ & \repo{experiments/ridgeline\_master.py}, \repo{experiments/ridgeline\_vertex.py} \\
the near-onset layer & \repo{experiments/boundary\_layer.py} \\
\bottomrule
\end{tabular}
\caption{Where each claim is checked. The banked notes carrying the full
records are \repo{results/onset-defect-law.md},
\repo{results/onset-defect-depth1-closed.md},
\repo{results/onset-defect-depths234.md} and
\repo{results/ridgeline-depth-amplitudes.md}.}
\label{tab:repro}
\end{table}

% The only citation here is the singularity-analysis machinery of
% \S\ref{sec:depth1}. Its mathematical antecedents are the companion papers,
% named in prose as they are throughout this directory, and its priority pass
% (\S\ref{sec:novelty}) found nothing else to cite.
\bibliographystyle{plain}
\bibliography{shared/refs}

\end{document}
